%----------------------------------------------------------------------------------------
%   CLASS DEFINITIONS (EMBEDDED RESUME.CLS)
%----------------------------------------------------------------------------------------
\documentclass[9pt,letterpaper]{article} % Font size and paper type

\usepackage[parfill]{parskip} % Remove paragraph indentation
\usepackage{array} % Required for boldface (\bf and \bfseries) tabular columns
\usepackage{ifthen} % Required for ifthenelse statements
\usepackage{charter}
\usepackage{enumitem}
\usepackage{xcolor} 
\usepackage[colorlinks=true, urlcolor=blue]{hyperref}
\usepackage[left=0.5in,top=0.25in,right=0.4in,bottom=0.2in, columnsep=0.5cm]{geometry} 
\input{glyphtounicode}
\usepackage[none]{hyphenat}
\usepackage{multicol}

\pdfgentounicode=1
\pagestyle{empty} % Suppress page numbers

\makeatletter % Begin internal command redefinitions

%----------------------------------------------------------------------------------------
%	HEADINGS COMMANDS: Commands for printing name and address
%----------------------------------------------------------------------------------------

\def \name#1{\def\@name{#1}} % Defines the \name command to set name
\def \@name {} % Sets \@name to empty by default

\def \addressSep {$\diamond$} % Set default address separator to a diamond

% One, two or three address lines can be specified 
\let \@addressone \relax
\let \@addresstwo \relax
\let \@addressthree \relax

% \address command can be used to set the first, second, and third address (last 2 optional)
\def \address #1{
  \@ifundefined{@addresstwo}{
    \def \@addresstwo {#1}
  }{
  \@ifundefined{@addressthree}{
  \def \@addressthree {#1}
  }{
     \def \@addressone {#1}
  }}
}

% \printaddress is used to style an address line (given as input)
\def \printaddress #1{
  \begingroup
    \def \\ {\addressSep\ }
    \centerline{#1}
  \endgroup
  \par
}

% \printname is used to print the name as a page header
\def \printname {
  \begingroup
    \hfil{{\namesize\bf \@name}}\hfil
    \nameskip\break
  \endgroup
}

%----------------------------------------------------------------------------------------
%	PRINT THE HEADING LINES
%----------------------------------------------------------------------------------------

\let\ori@document=\document
\renewcommand{\document}{
  \ori@document  % Begin document
  \printname % Print the name specified with \name
  \@ifundefined{@addressone}{}{ % Print the first address if specified
    \printaddress{\@addressone}}
  \@ifundefined{@addresstwo}{}{ % Print the second address if specified
    \printaddress{\@addresstwo}}
     \@ifundefined{@addressthree}{}{ % Print the third address if specified
    \printaddress{\@addressthree}}
}

%----------------------------------------------------------------------------------------
%	SECTION FORMATTING
%----------------------------------------------------------------------------------------

% Defines the rSection environment for the large sections within the CV
\newenvironment{rSection}[1]{ % 1 input argument - section name
  \sectionskip
 \textbf{\Large #1} % Section title
  \sectionlineskip
  \hrule % Horizontal line
  \begin{list}{}{ % List for each individual item in the section
    \setlength{\leftmargin}{-0em} % Margin within the section
  }
  \item[]
}{
  \end{list}
}

%----------------------------------------------------------------------------------------
%	WORK EXPERIENCE FORMATTING
%----------------------------------------------------------------------------------------

\newenvironment{rSubsection}[4]{ % 4 input arguments - company name, year(s) employed, job title and location
 {\bf #1} \hfill {#2} % Bold company name and date on the right
 \ifthenelse{\equal{#3}{}}{}{ % If the third argument is not specified, don't print the job title and location line
  \\
  {\em #3} \hfill {\em #4} % Italic job title and location
  }\smallskip
  \begin{list}{$\cdot$}{\leftmargin=0em} % \cdot used for bullets, no indentation
   \itemsep -0.5em \vspace{-0.5em} % Compress items in list together for aesthetics
  }{
  \end{list}
  \vspace{0em} % Some space after the list of bullet points
}

% The below commands define the whitespace after certain things in the document - they can be \smallskip, \medskip or \bigskip
\def\namesize{\huge} % Size of the name at the top of the document
\def\addressskip{} % The space between the two address (or phone/email) lines
\def\sectionlineskip{\smallskip} % The space above the horizontal line for each section 
\def\nameskip{\vspace{-0.5em}} % The space after your name at the top
\def\sectionskip{} % The space after the heading section

\makeatother % End internal command redefinitions

%----------------------------------------------------------------------------------------
%   DOCUMENT START
%----------------------------------------------------------------------------------------
\name{[First Name] [Last Name], ([Credentials]) \vspace{.2em} \hrule} % Your name

\address{[City, State, Country] \textbullet\ [Email Address] \textbullet\ {\href{[Website URL]}{[Website URL]}} \textbullet\ {\href{[LinkedIn Profile URL]}{linkedin.com/in/[profile]}} }

\begin{document}
\textbf{[Professional Summary: A brief, high-impact statement summarizing your experience, key achievements, and core value proposition. Mention years of experience, key industries, and top skills. Example: Results-oriented Software Engineer with 5+ years of experience in...]}


\begin{rSection}{Education}
\begin{rSubsection}{[University Name]}{[Location]}{\textbf{[Degree Title]}}{[Start Date] - [End Date]}
\item \textbf{[Minor / Specialization]} \hfill [Date]
\end{rSubsection}

\begin{rSubsection}{[University Name]}{[Location]}{\textbf{[Degree Title]}}{[Start Date] - [End Date]}
\item[]
\end{rSubsection}
\end{rSection}

%----------------------------------------------------------------------------------------
%   WORK EXPERIENCE SECTION
%----------------------------------------------------------------------------------------

\begin{rSection}{Work Experience}
\begin{rSubsection}
{[Job Title], \href{[Company URL]}{[Company Name]} - [Location]}{[Start Date] - [End Date]}{}{}
\begin{itemize}[left=-1.5em, labelsep=.1cm, labelwidth=.2cm, itemsep=0.0em]
\item Accomplished \textbf{[Result X]} as measured by \textbf{[Metric Y]}, by \textbf{[Action Z]}.
\item Accomplished \textbf{[Result X]} as measured by \textbf{[Metric Y]}, by \textbf{[Action Z]}.
\item Accomplished \textbf{[Result X]} as measured by \textbf{[Metric Y]}, by \textbf{[Action Z]}.
\end{itemize}
\end{rSubsection}

\begin{rSubsection}
{[Job Title], \href{[Company URL]}{\textit{[Company Name]}} - [Location]}{[Start Date] - [End Date]}{}{}
\begin{itemize}[left=-1.5em, labelsep=.1cm, labelwidth=.2cm, itemsep=0.0em]
\item Accomplished \textbf{[Result X]} as measured by \textbf{[Metric Y]}, by \textbf{[Action Z]}.
\item Accomplished \textbf{[Result X]} as measured by \textbf{[Metric Y]}, by \textbf{[Action Z]}.
\item Accomplished \textbf{[Result X]} as measured by \textbf{[Metric Y]}, by \textbf{[Action Z]}.
\end{itemize}
\end{rSubsection}

\end{rSection}

\begin{rSection}{Projects (Product Work) {\fontsize{8}{9}\selectfont (Details in portfolio.)}}
\begin{rSubsection}{\href{[Project URL]}{[Project Name]}}{}{}{}
\item[] \vspace{-0.5em}
\begin{itemize}[left=-1.5em, labelsep=.1cm, labelwidth=.2cm, itemsep=0.0em]
\item Accomplished \textbf{[Result X]} as measured by \textbf{[Metric Y]}, by \textbf{[Action Z]}.
\item Accomplished \textbf{[Result X]} as measured by \textbf{[Metric Y]}, by \textbf{[Action Z]}.
\end{itemize}
\end{rSubsection}

\begin{rSubsection}{\href{[Project URL]}{[Project Name]}}{}{}{}
\item[] \vspace{-0.5em}
\begin{itemize}[left=-1.5em, labelsep=.1cm, labelwidth=.2cm, itemsep=0.0em]
\item Accomplished \textbf{[Result X]} as measured by \textbf{[Metric Y]}, by \textbf{[Action Z]}.
\item Accomplished \textbf{[Result X]} as measured by \textbf{[Metric Y]}, by \textbf{[Action Z]}.
\end{itemize}
\end{rSubsection}
\end{rSection}

\begin{rSection}{Skills and Tools}

\begin{itemize}[left=-1.5em, labelsep=.1cm, labelwidth=.2cm, itemsep=0.0em]
    \item \textbf{[Category 1]:} [Skill 1], [Skill 2], [Skill 3] | \textbf{[Sub-Category]}: [Skill], [Skill] | \textbf{[Platform]}: [Tool], [Tool]
    \item \textbf{[Category 2]:} [Skill 1], [Skill 2], [Skill 3] | \textbf{[Sub-Category]}: [Skill], [Skill] | \textbf{[Platform]}: [Tool], [Tool]
    \item \textbf{[Category 3]:} \textbf{[Attribute]}: [Trait], [Trait] | \textbf{[Strategy]}: [Strategy 1], [Strategy 2] | \textbf{[Management]}: [Mgmt Skill 1], [Mgmt Skill 2]
    \item \textbf{[Category 4]:} Methodologies: [Agile/Scrum] | \textbf{Languages}: [Lang 1], [Lang 2] | \textbf{Technologies}: [Tech 1], [Tech 2] | \textbf{DevOps}: [Tool 1], [Tool 2]
\end{itemize}

\end{rSection}


\begin{rSection}{Certification}
\item[] \vspace{-1em}
\begin{itemize}[left=-1.5em, labelwidth=.1cm, itemsep=0.0em]
\item \href{[Certificate URL]}{[Certificate Name]}
\item \href{[Certificate URL]}{[Certificate Name]}
\item \href{[Certificate URL]}{[Certificate Name]}
\item \href{[Certificate URL]}{[Certificate Name]}
\end{itemize}
\end{rSection}

\end{document}
